\documentclass[11pt]{article}

\usepackage[margin=1in]{geometry}
\usepackage{amsmath,amssymb,mathtools}
\usepackage{booktabs}
\usepackage{enumitem}
\usepackage{hyperref}
\usepackage[T1]{fontenc}
\usepackage[utf8]{inputenc}
\usepackage{lmodern}
\usepackage{listings}

\setlist[itemize]{nosep,leftmargin=2em}
\setlist[enumerate]{nosep,leftmargin=2em}
\lstset{basicstyle=\ttfamily\small,breaklines=true,frame=single}

\newcommand{\faultdim}{d_\Gamma}
\newcommand{\R}{\mathbb{R}}
\newcommand{\Gammah}{\Gamma_h}
\newcommand{\wh}{\omega}

\title{Phase-Field-Aware Reconstructed Fault\\Technical Specification for ASPECT/deal.II}
\author{}
\date{Revision 3}

\begin{document}
\maketitle

\section*{Status of this document}
This document and \texttt{current\_design.md} are the implementation sources of
truth for the reconstructed-fault framework. They are maintained as a
consistent pair. Older redesign documents and derivation notes are background
only unless an approved equation is first promoted into both authoritative
documents. Codex should implement only what is specified here and should
report ambiguities or conflicts instead of silently choosing a different
algorithm.

\section*{Reader map}
The framework has four layers:
\begin{enumerate}
  \item \textbf{Replicated fault registry}: ordered 2-D polyline geometry,
        connectivity, generic committed constitutive properties, and the
        distinguished slip-rate field.
  \item \textbf{Fault Q1 operations}: segment-local interpolation and
        fault-sized mass/projection systems on the ordered polyline, without a
        duplicate runtime triangulation.
  \item \textbf{Coupling operators and caches}: cached particle-to-fault and fault-to-quadrature-point mappings, plus the reduced projection operators used inside nonlinear iterations.
  \item \textbf{Distributed bulk fields}: the ASPECT mesh, phase field, particles, and Stokes quadrature points remain distributed over MPI ranks.
\end{enumerate}

The central data-flow pattern is
\[
\boxed{
\text{distributed bulk data}
\;\longrightarrow\;
\text{local assembly/evaluation}
\;\longrightarrow\;
\text{MPI reduction when needed}
\;\longrightarrow\;
\text{small replicated fault problem}
}
\]
not ``gather all bulk data on rank 0.''


\section{Existing implementation interface}
\label{sec:existing-interface}

\subsection{Reference policy}
The existing code is divided below into three categories:
\begin{itemize}
  \item \textbf{MUST reuse}: an existing API or data structure is authoritative and a duplicate implementation/parameter shall not be introduced;
  \item \textbf{MAY reuse/refactor}: an isolated implementation is useful, but the surrounding old algorithm is not authoritative;
  \item \textbf{MUST NOT depend on}: the implementation belongs to the abandoned core-phase-field/original phase-field-RSF algorithm and shall not become a dependency of the reconstructed-fault framework.
\end{itemize}

\subsection{Phase-field module}
The authoritative implementation is
\begin{center}
\texttt{include/aspect/phase\_field.h},\qquad
\texttt{source/simulator/phase\_field.cc}.
\end{center}
The existing \texttt{PhaseFieldHandler<dim>} already declares the bulk \texttt{phase\_field} variable with \texttt{FE\_Q<dim>(1)} and owns the machinery used to evolve the distributed phase field.  The reconstructed-fault implementation shall use this existing FE variable; it shall not introduce a second phase-field representation.

\paragraph{Phase-field ranges and thresholds.}
The authoritative physical phase-field bounds are provided by
\begin{center}
\texttt{MaterialModel::PhaseFieldModel<dim>::get\_phase\_field\_range()}.
\end{center}
This range is exactly $[0,1]$. The reconstruction
activation threshold and the fault-model upper admissibility threshold are
separate quantities supplied by
\texttt{get\_phase\_field\_activation\_threshold()} and
\texttt{get\_phase\_field\_upper\_admissibility\_threshold()}.
Reconstruction and phase-field mesh refinement shall continue to use the
activation threshold. Prescribed fault-core validation and legacy slip-rate
normalization shall continue to use both model-specific admissibility
thresholds. No independent duplicate threshold shall be introduced.
The generic upper-admissibility default is the physical upper bound $1$. A
model that uses an algorithm requiring a strict interior endpoint must
override it; in particular, \texttt{PhaseFieldFault} supplies $0.99$ for the
legacy slip-rate normalizer.

\paragraph{Length scale.}
The phase-field length scale $\ell$ is already parsed by \texttt{PhaseFieldHandler<dim>} and stored in its \texttt{PhaseField::GeometricFunction}; \texttt{GeometricFunction::get\_length\_scale()} already returns it.  \textbf{MUST reuse:} the reconstructed-fault module shall not define a duplicate fracture length scale.  Because \texttt{PhaseFieldHandler<dim>} currently does not expose $\ell$ directly, the preferred change is to add a thin const accessor such as
\begin{center}
\texttt{PhaseFieldHandler<dim>::get\_length\_scale() const}
\end{center}
that forwards to \texttt{geometric\_function->get\_length\_scale()}.

\paragraph{Associated particle infrastructure.}
\texttt{PhaseFieldHandler<dim>} already exposes
\begin{center}
\texttt{get\_associated\_particle\_manager()},\qquad
\texttt{get\_grid\_cache()}.
\end{center}
These existing access points should be reused where appropriate.  In particular, the reconstructed-fault coupling should obtain the particle manager associated with the phase-field calculation through the existing handler instead of discovering a second particle set independently.

\paragraph{Core-phase-field functionality.}
\textbf{MUST NOT depend on:} the reconstructed-fault framework shall not call or require the old core-phase-field machinery, including \texttt{PhaseFieldHandler<dim>::extend\_core\_phase\_field()}, the \texttt{pre\_extend\_core\_phase\_field} signal, or algorithms whose inputs fundamentally require \texttt{core\_phase\_field}.  This includes the old uses of \texttt{slip\_rate\_localization\_factor(..., core\_phase\_field)} and \texttt{stationary\_crack\_driving\_force(..., core\_phase\_field)} in the failed RSF localization algorithm.  The existing core-phase-field FE variable may remain temporarily for backward compatibility with old plugins, but the new reconstructed-fault implementation shall neither read it nor require it.

\paragraph{Old particle interpolation helper.}
\texttt{PhaseFieldUtilities::interpolate\_from\_particles\_in\_crack\_zone()} is tied to interpolation of particle properties in the old diffuse crack-zone representation.  It shall not be used as the particle-to-sharp-fault projection operator defined in this specification.  Isolated implementation techniques from it may be consulted, but the new operator must assemble directly onto the reconstructed ordered-polyline Q1 space.

\subsection{Particle-domain module}
The authoritative implementation is
\begin{center}
\texttt{include/aspect/particle/particle\_domain.h},\qquad
\texttt{source/particle/particle\_domain.cc}.
\end{center}
The existing \texttt{Particle::ParticleDomainHandler<dim>} and \texttt{ParticleDomainAccessor<dim>} provide particle-domain geometry and CPDI information.  In particular,
\begin{center}
\texttt{get\_particle\_domain(particle\_index).volume()}
\end{center}
returns the physical measure of the particle domain.

\textbf{MUST reuse:} this particle-domain volume is the preferred sampling measure $m_p$ for particle-to-fault projection.  Do not approximate it by $|K|/N_{p,K}$ when the particle-domain volume is available.

The module also provides particle-domain face/neighbour data and CPDI relevant-vertex weighting functions.  These may be reused for tasks for which they were designed, but the CPDI weighting functions are bulk-grid transfer weights and are \emph{not} the shape functions of the reconstructed fault.  Therefore \texttt{weighting\_function\_value()} and \texttt{weighting\_function\_gradient()} shall not replace the fault Q1 basis $N_a(\boldsymbol\xi_\Gamma)$ in the particle-to-fault projection.

\subsection{Fault friction rheology}
The authoritative implementation is
\begin{center}
\texttt{include/aspect/material\_model/rheology/fault\_friction.h},\qquad
\texttt{source/material\_model/rheology/fault\_friction.cc}.
\end{center}
The class is
\begin{center}
\texttt{MaterialModel::Rheology::FaultFriction<dim>}.
\end{center}
It selects either the existing stateful rate-and-state law or the stateless
rate-dependent law
\begin{equation}
  \mu(V)=\mu_d+\frac{\mu_s-\mu_d}{1+V/V_c},
  \qquad
  \frac{\partial\mu}{\partial V}
  =-\frac{(\mu_s-\mu_d)V_c}{(V_c+V)^2}.
\end{equation}
For a material mixture, first arithmetic-average each parameter,
\begin{equation}
  \bar\mu_s=\sum_m f_{\Gamma,m}\mu_{s,m},\qquad
  \bar\mu_d=\sum_m f_{\Gamma,m}\mu_{d,m},\qquad
  \bar V_c=\sum_m f_{\Gamma,m}V_{c,m},
\end{equation}
and then evaluate one rate-dependent law with the averaged parameters.  Do not
average separately evaluated nonlinear friction responses.

\textbf{MUST reuse:} the reconstructed-fault constitutive model shall use the existing RSF implementation for at least
\begin{itemize}
  \item \texttt{update\_state(V, theta\_old, dt)};
  \item \texttt{friction\_coefficient(volume\_fractions, V, theta)};
  \item \texttt{friction\_coefficient\_derivative\_wrt\_slip\_rate(...)};
  \item \texttt{compute\_time\_step(...)};
  \item existing getters for $V_0$, $V_{\min}$, and $D_c$ where required;
  \item the existing parameter declaration/parsing for $V_0$, $V_{\min}$, $D_c$, $\mu_0$, $a$, $b$, and the regularization option.
\end{itemize}
The new reconstructed-fault model shall not duplicate these formulas or define an independent set of RSF parameters.

The existing \texttt{Reference friction coefficients} parameter represents
$\mu_0$ for rate-and-state friction and $\mu_s$ for rate-dependent friction.
Rate-dependent friction additionally declares composition-aware
\texttt{Dynamic friction coefficients}, $\mu_d$, with default $0.4$, and
\texttt{Characteristic weakening slip rates}, $V_c$, with default
$10^{-6}\,\si{\meter\per\second}$.  Componentwise admissibility requires
$V_c>0$ and $0\leq\mu_d\leq\mu_s$.

The stateful coefficient and derivative interfaces take $\theta$.  Separate
stateless overloads omit $\theta$; using an overload that does not match the
selected law is an API error.  Both laws require $V\geq V_{\min}>0$.  There is
no constitutive upper slip-rate clamp; an admissible trial value is evaluated
as supplied.  The rate-dependent law owns no state and imposes no friction
timestep restriction, so \texttt{compute\_time\_step(...)} returns the largest
finite \texttt{double}.  Stage E is constitutive-only: it does not add
persistent $\Theta$, a surface residual, mechanical coupling, or nonlinear
commit/rollback hooks.

\subsection{Original phase-field-RSF material model: reference only}
The file
\begin{center}
\texttt{tmp/material\_model/phase\_field\_rsf.cc}
\end{center}
implements the abandoned diffuse/core-phase-field RSF architecture.  \textbf{MUST NOT use as an architectural base class or data-ownership model.}  In particular, the new implementation shall not copy its dependencies on the core phase field, particle-stored normal/slip directions, or the old diffuse slip-localization factor.

\textbf{MAY reuse/refactor:} selected isolated numerical procedures may be extracted after verifying that their inputs remain valid in the sharp-fault formulation.  The most important example is the local nonlinear return-mapping/root-solving procedure currently contained in \texttt{PhaseFieldRSF<dim>::perform\_return\_mapping()}.  Its scalar RSF root equation, derivative evaluation, convergence checks, failure reporting, and use of \texttt{RateStateFriction<dim>} are useful references.  They should be refactored into a geometry-independent helper operating on explicit local inputs rather than copying the complete particle loop.

The current signal timing is also a useful reference: the old model performs return mapping from the \texttt{pre\_assemble\_stokes\_system} signal and updates history after the nonlinear solver.  The new implementation may preserve compatible timing where appropriate, but it shall not reuse the \texttt{pre\_extend\_core\_phase\_field} path or \texttt{update\_particles\_in\_emerging\_crack\_zone()}.

\subsection{Original phase-field-RSF particle property: do not use as state layout}
The file
\begin{center}
\texttt{tmp/particle\_property/phase\_field\_rsf.cc}
\end{center}
is specific to the failed diffuse-fault algorithm.  It currently defines particle properties including \texttt{slip\_rate}, \texttt{slip\_state}, \texttt{normal\_direction}, and \texttt{slip\_direction}.  \textbf{MUST NOT use this layout as the authoritative state layout for the reconstructed fault.}  In the new formulation, fault state and fault geometry live on the reconstructed codimension-one fault according to this specification.

Individual utility code may be consulted or refactored when it is independent of the abandoned geometry/state assumptions, but no function from this plugin is reused merely because it already has a similar name.  In particular, its old crack-zone particle interpolation and emerging-crack normal/slip-direction initialization are not part of the reconstructed-fault algorithm.

\subsection{Rule for future interface additions}
When the reconstructed-fault code needs information that already exists privately in one of the reusable modules, prefer adding a narrow const accessor to the authoritative module over duplicating the parameter or recomputing the quantity elsewhere.  Such additions must expose existing state, not create a second source of truth.

\section{Scope and non-goals}

The public interface shall be dimension-independent:
\[
\faultdim = \texttt{dim}-1.
\]
Thus \texttt{dim=2} corresponds to a one-dimensional fault curve embedded in two dimensions, and a future \texttt{dim=3} instantiation corresponds to a two-dimensional fault surface embedded in three dimensions.

The current implementation only needs a complete propagation algorithm for \texttt{dim=2}.  When the 2-D and 3-D implementations differ substantially, only the 2-D implementation is required now.  The generic interface must nevertheless avoid assumptions that would prevent a future 3-D implementation.

The current implementation shall not attempt:
\begin{itemize}
  \item 3-D crack-front propagation or 3-D surface remeshing;
  \item branching or nucleation away from an existing fault;
  \item merging of independent faults;
  \item topology changes;
  \item a global Cartesian graph representation such as $y=f(x)$ or $z=f(x,y)$;
  \item a global $p_1\times p_2$ structural parameterization;
  \item distributed ownership of the fault mesh.
\end{itemize}

\section{Fundamental invariants}

\subsection{Codimension-one manifold}
The fault is represented as a codimension-one manifold embedded in the bulk domain.  The public interface shall operate on fault cells, physical points, reference coordinates on fault cells, and fault normals.  It shall not expose structured-grid indices such as \texttt{i}, \texttt{j}, \texttt{p1}, or \texttt{p2}.

\subsection{Append-only persistent geometry}
Once a fault vertex is committed to the persistent fault, its physical location never changes.  Geometry correction and smoothing are performed only on temporary candidate geometry before it is committed.

This rule prevents repeated remapping of rate-and-state history variables.

\subsection{Replicated fault, distributed bulk}
Every MPI rank stores the complete persistent fault registry. The ASPECT bulk
triangulation, FE fields, and particles remain distributed. No second runtime
fault mesh duplicates the ordered-polyline geometry in the current 2-D
implementation.

\subsection{No abuse of triangulation user data}
Persistent fault geometry and constitutive state shall \emph{not} be stored in
\texttt{user\_data}, \texttt{user\_index}, or \texttt{user\_flag}. These
deal.II mechanisms may be used only for short-lived algorithmic bookkeeping
when appropriate.

\section{Geometry and state representation}

\subsection{Persistent registry}
The persistent application-owned registry is the lifetime-stable
representation of fault geometry and state. In the current fixed-geometry 2-D
implementation, one fault is an ordered sequence of points, consecutive points
define segments implicitly, and a geometry-version counter invalidates
reconstructible caches.
\begin{lstlisting}[language=C++]
template <int dim>
class ReconstructedFault
{
  std::vector<dealii::Point<dim>> vertices;
  std::vector<double> generic_vertex_properties;
  std::uint64_t geometry_version;
};
\end{lstlisting}

The manager defines one shared runtime property schema before geometry exists.
Every fault stores one contiguous vertex-major property array. This pool
contains committed physical/material state only. Constitutive meaning belongs
to the material model; temporary and trial data do not belong in the pool.

\subsection{Surface Q1 representation}
The ordered segment endpoints are the Q1 nodes. Segment interpolation uses
\(N_0=1-\xi\), \(N_1=\xi\). Fault-sized consistent mass matrices and their
tridiagonal factors provide projection. A future 3-D surface representation is
deferred and must not force duplicate geometry into the current implementation.

\subsection{Surface location}
A generic surface location is
\[
\boxed{
\text{SurfaceLocation}
=
(\text{fault-cell identifier},\,\boldsymbol\xi_\Gamma)
}
\]
with
\[
\boldsymbol\xi_\Gamma\in[0,1]^{\faultdim}.
\]
For \texttt{dim=2}, $\boldsymbol\xi_\Gamma=(\xi)$.  A future \texttt{dim=3} implementation would use $(\xi,\eta)$.

The public fault interface should provide operations conceptually equivalent to
\begin{lstlisting}
closest_point(x)
position(surface_location)
normal(surface_location)
interpolate_surface_field(surface_location, field)
geometry_version()
\end{lstlisting}

\section{Phase-field-aware reconstruction core}

Define a normalized crack variable
\[
d_h(\mathbf{x})\in[0,1],
\]
with $d_h=0$ in intact material and $d_h=1$ at the crack core.  If the FE phase-field convention is opposite, convert it before reconstruction.

Let \(\phi_{\rm act}\) be supplied by
\texttt{get\_phase\_field\_activation\_threshold()} and let
\(\phi_{\rm adm}\) be supplied by
\texttt{get\_phase\_field\_upper\_admissibility\_threshold()}.
The physical range returned by
\texttt{get\_phase\_field\_range()} remains exactly \([0,1]\).
Reconstruction support uses \(\phi_{\rm act}\), not the physical lower bound.
Use
\[
\wh(d)
=
\begin{cases}
0, & d\le \phi_{\rm act},\\[1mm]
\left[
\dfrac{\min(d,\phi_{\rm adm})-\phi_{\rm act}}
{\phi_{\rm adm}-\phi_{\rm act}}
\right]^p,
& d>\phi_{\rm act},
\end{cases}
\]
with initial exponent $p=2$.

For a provisional fault point $\mathbf{R}$ with unit normal $\mathbf{n}$, sample
\[
\mathbf{x}(\eta)=\mathbf{R}+\eta\mathbf{n},
\qquad
-R_n\le\eta\le R_n.
\]
At sample points $\eta_k$, evaluate
\[
d_k=d_h(\mathbf{R}+\eta_k\mathbf{n}),
\qquad
w_k=q_k\wh(d_k).
\]
Then compute
\[
M=\sum_k w_k,
\]
\[
\boxed{
\delta=\frac{\sum_k w_k\eta_k}{\sum_k w_k}
},
\]
\[
\boxed{
\sigma^2
=
\frac{\sum_k w_k(\eta_k-\delta)^2}{M}
},
\]
and
\[
d_{\max}=\max_k d_k.
\]

The phase-field-centered point is
\[
\boxed{
\mathbf{Q}=\mathbf{R}+\delta\mathbf{n}.
}
\]

A suitable initial transverse radius and sampling spacing are
\[
R_n=\max(2\ell,2h_{\mathrm{pf}}),
\qquad
\Delta\eta\lesssim0.5h_{\mathrm{pf}}.
\]

For mature fault locations define
\[
M_{\mathrm{ref}}=\operatorname{median}_{i\in\mathcal I_{\mathrm{mature}}}M_i,
\qquad
\rho_i=\frac{M_i}{M_{\mathrm{ref}}}.
\]

\section{Main reconstruction does not need a bulk-node-to-fault map}

The normal-line samples and 2-D forward-PCA samples are arbitrary physical points generated from the current reconstructed fault.  Their phase-field values are obtained directly from the distributed FE solution.

Therefore the main propagation algorithm shall \emph{not} build or maintain a mapping
\[
\text{bulk FE node}\longrightarrow\text{fault cell}.
\]

A Xu-style bulk-node point-cloud regression is retained only as an optional initialization/fallback algorithm for a future case in which no initial sharp fault geometry is prescribed.  It is not part of the normal propagation path.

\section{2-D propagation algorithm}

This section is specific to \texttt{dim=2}.

The persistent fault is an ordered chain
\[
\mathbf{X}_0,\mathbf{X}_1,\ldots,\mathbf{X}_{N-1}.
\]
For an interior vertex,
\[
\mathbf{t}_i
=
\frac{\mathbf{X}_{i+1}-\mathbf{X}_{i-1}}
{|\mathbf{X}_{i+1}-\mathbf{X}_{i-1}|},
\qquad
\mathbf{n}_i=(-t_{i,y},t_{i,x}).
\]
At an endpoint, use the adjacent segment.

Let the nominal structural spacing satisfy approximately
\[
0.5h_{\mathrm{pf}}\lesssim h_\Gamma\lesssim h_{\mathrm{pf}}.
\]

For a propagating tip $\mathbf{X}_T$ with tangent $\mathbf{t}_T$ and normal $\mathbf{n}_T$, define a forward search patch
\[
\Omega_T
=
\left\{
\mathbf{x}:0<(\mathbf{x}-\mathbf{X}_T)\cdot\mathbf{t}_T<R_T,
\quad
|(\mathbf{x}-\mathbf{X}_T)\cdot\mathbf{n}_T|<R_N
\right\},
\]
with initial scales $R_T\sim2\ell$--$3\ell$ and $R_N\sim2\ell$.

At sample points $\mathbf{x}_a$ use weights
\[
w_a=q_a\wh(d_h(\mathbf{x}_a)).
\]
Compute the weighted centroid and covariance
\[
\boldsymbol\mu
=
\frac{\sum_a w_a\mathbf{x}_a}{\sum_a w_a},
\]
\[
\mathbf{C}_T
=
\frac{
\sum_a w_a(\mathbf{x}_a-\boldsymbol\mu)\otimes(\mathbf{x}_a-\boldsymbol\mu)
}{\sum_a w_a}.
\]
If $\lambda_1\ge\lambda_2$ are its eigenvalues and $\mathbf{v}_1$ corresponds to $\lambda_1$, orient
\[
\mathbf{t}_{\mathrm{PCA}}
=
\operatorname{sign}(\mathbf{v}_1\cdot\mathbf{t}_T)\mathbf{v}_1.
\]
Define
\[
A_T
=
\frac{\lambda_1-\lambda_2}{\lambda_1+\lambda_2+\varepsilon},
\qquad
\gamma=\min\left(1,\frac{A_T}{A_0}\right),
\]
and
\[
\mathbf{t}_{\mathrm{pred}}
=
\operatorname{normalize}
\left[(1-\gamma)\mathbf{t}_T+\gamma\mathbf{t}_{\mathrm{PCA}}\right].
\]
Enforce a configurable maximum turning angle, initially about $45^\circ$.

Predict
\[
\mathbf{R}_{\mathrm{new}}
=
\mathbf{X}_T+h_\Gamma\mathbf{t}_{\mathrm{pred}},
\]
then apply the transverse phase-field centering operation to obtain
\[
\mathbf{Q}_{\mathrm{new}}
=
\mathbf{R}_{\mathrm{new}}+\delta_{\mathrm{new}}\mathbf{n}_{\mathrm{pred}}.
\]

A candidate must satisfy configurable criteria based on
\[
\rho_{\mathrm{new}},\qquad d_{\max,\mathrm{new}},\qquad
|\delta_{\mathrm{new}}|,\qquad\sigma_{\mathrm{new}},\qquad
\Delta\theta.
\]
It must also avoid intersection with nonadjacent persistent fault cells.

Accepted candidates first form a temporary extension.  Only after geometric regularization and a final quality check are they committed to the persistent registry.

\section{Temporary-geometry regularization}

Only temporary candidate vertices may move.  Write
\[
\mathbf{X}_i^{\mathrm{new}}
=
\mathbf{X}_i+u_i\mathbf{n}_i.
\]
The phase-field profile suggests $u_i\approx\delta_i$.

Use the generic objective
\[
J(\mathbf{u})
=
\frac12(\mathbf{u}-\boldsymbol\delta)^T\mathbf{C}
(\mathbf{u}-\boldsymbol\delta)
+
\frac{\lambda_s}{2}\|L_\Gamma\mathbf{u}\|^2,
\]
where $\mathbf{C}=\operatorname{diag}(c_i)$ contains confidence weights.

For the current 2-D implementation,
\[
(L_\Gamma\mathbf{u})_i
=u_{i-1}-2u_i+u_{i+1}.
\]
Persistent anchor vertices satisfy $u_i=0$.  Solve
\[
\boxed{
(\mathbf{C}+\lambda_sL_\Gamma^TL_\Gamma)\mathbf{u}
=
\mathbf{C}\boldsymbol\delta.
}
\]

This problem contains only temporary fault unknowns and is therefore a small replicated/serial problem.  It is \emph{not} assembled from millions of bulk nodes and does not require a distributed linear solver.

Repeat normal evaluation, phase-field sampling, and this correction for at most a few geometry iterations, initially $N_{\mathrm{geom,max}}=3$.

\section{Parallel coupling operators}

There are three distinct mappings in the coupled model:
\[
P_{\Gamma\leftarrow h}:
\text{distributed FE phase field}\rightarrow\text{fault geometry},
\]
\[
P_{\Gamma\leftarrow p}:
\text{distributed particles}\rightarrow\text{fault surface fields},
\]
\[
P_{q\leftarrow\Gamma}:
\text{replicated fault fields}\rightarrow\text{local Stokes quadrature points}.
\]

They have different MPI patterns:

\begin{center}
\begin{tabular}{llll}
\toprule
Operator & Source & Destination & Parallel pattern \\
\midrule
$P_{\Gamma\leftarrow h}$ & distributed FE field & replicated geometry & distributed point evaluation \\
$P_{\Gamma\leftarrow p}$ & distributed particles & replicated surface field & local assembly + MPI reduction + replicated solve \\
$P_{q\leftarrow\Gamma}$ & replicated surface field & local QPs & cached local interpolation only \\
\bottomrule
\end{tabular}
\end{center}

\section{Distributed phase-field evaluation}

The normal profiles and forward-PCA patches contain arbitrary physical sample points.  These values shall be obtained with deal.II distributed point-evaluation infrastructure, conceptually \texttt{Utilities::MPI::RemotePointEvaluation} or the corresponding ASPECT wrapper.

For adaptive normalization profiles, deterministic fault-major, segment-major,
quadrature-point IDs shall be divided into balanced contiguous owner ranges.
Each owner generates and updates only its profiles, while all ranks participate
in every batched point-evaluation collective. There is no rank-zero profile
coordinator.

No rank shall gather the complete bulk FE field or a complete bulk-node point cloud.

\subsection{Normalization-profile material mixture}
For every reconstructed-fault quadrature point determine one surface material
mixture $\mathbf f_\Gamma$ by projecting chemical compositional fields from
particles to the Q1 fault surface and converting the interpolated surface
composition to material fractions. Hold this mixture fixed for the complete
$+\mathbf n/-\mathbf n$ profile and evaluate
\[
\bar g(\phi,\mathbf f_\Gamma)=\sum_m f_{\Gamma,m}g_m(\phi),\qquad
I_h=\int\left(1/\bar g-1\right)d\zeta.
\]
Do not evaluate a different composition mixture at individual normal-profile
samples. This model requires material transitions to have characteristic width
$L_{\rm mat}\gg\ell$; transitions on the diffuse-fault scale or narrower are
outside its applicability.

Register every chemical compositional field as a separate scalar generic fault
property named
\texttt{phase field fault chemical composition <field-name>}, preserving the
chemical-field introspection order. This keeps the corresponding visualization
arrays individually named. If no chemical compositional fields exist, register
no chemical fault property and interpret an empty composition vector as the
background-only material mixture.

The physical lower bound is zero, but the unconstrained Q1 phase-field solve
may produce a small negative numerical undershoot. For $I_h$ only, the
degradation function shall therefore be evaluated using
\[
\phi_{\mathrm{eff}}=\max(\phi_h,0).
\]
The activation threshold shall not be used as this clamp. The minimum raw
sampled phase field shall be tracked across all profiles and MPI ranks. The
internal empirical dimensionless error-detection threshold is $10^{-4}$; it
is not a physical parameter, solver tolerance, convergence criterion, or
user-adjustable numerical convergence parameter. A global minimum below
$-10^{-4}$ shall fail with diagnostics giving the raw minimum, the threshold,
and the sample location. Every raw sample must be finite. The upper end shall
not be clipped: a raw value above one is an
invariant failure, while the independent constitutive condition $\bar g>0$
remains authoritative. Thus $\bar g=0$, including a possible $\phi=1$
singularity, shall be reported explicitly as an $I_h$ singularity and shall
not be classified as a phase-field range violation. A tail value $\phi_h=0$
is admissible.

The consistent Q1 projection of positive quadrature-point normalization
integrals is retained. Consistent projection is not positivity preserving in
general, so positivity of its nodal coefficients is a tested property of the
smooth intended fault profiles rather than a generic mathematical invariant.

Verification shall include an analytic profile integral, quadrature- and
tail-tolerance convergence, and a distributed Q1 bulk-mesh refinement study
against an independently evaluated reference integral. A lifecycle smoke test
may use an artificially low reconstruction activation threshold solely to
provide geometry, but in that fixture reconstruction behavior is explicitly
outside the test's purpose.

\section{Common cohesive state}

\texttt{MaterialModel::PhaseFieldFault} owns the cohesive constitutive law.
The generic fault-property pool stores exactly the committed scalar traction
\(T^{\rm coh}_{k-1}\) and committed previous normalization
\(I_{h,k-1}\). The reconstructed-fault manager provides storage and
checkpointing only. Current \(I_{h,k}\), trial responses, and diagnostics
remain private material-model data.

With the Maxwell coefficients
\[
\beta_k=\exp(-\Delta tG/\eta),
\qquad
\kappa_k=-\eta\,\operatorname{expm1}(-\Delta tG/\eta),
\]
the non-committing surface update is
\[
\boxed{
T^{\rm coh}_k
=
\frac{\kappa_kV_k+
\beta_k I_{h,k-1}T^{\rm coh}_{k-1}}{I_{h,k}}.
}
\]
The exact pointwise diffuse crack-strain-rate magnitude is
\[
\boxed{
\upsilon_k
=
\frac{h_k}{I_{h,k}}V_k
+
\frac{\beta_kT^{\rm coh}_{k-1}}{\kappa_k}
\left[
h_k\frac{I_{h,k-1}}{I_{h,k}}-h_{k-1}
\right].
}
\]
The bracketed contribution is the history correction. It integrates to zero,
so
\[
\int\upsilon_k\,d\zeta=V_k.
\]
For current phase field and committed history frozen during a nonlinear
linearization,
\[
\frac{\partial\upsilon_k}{\partial V_k}
=\frac{h_k}{I_{h,k}}.
\]
The previous pointwise \(h_{k-1}\) is evaluated from ASPECT's previous FE
phase-field solution. It is not duplicated as fault history.

\subsection{Initial committed cohesive state}

For an initially reconstructed pre-existing fault, \(H_p\) and \(\phi_p\) are
particle-local quantities. Interpolate the
already projected Q1 chemical-composition fields at particle \(p\)'s cached
fault coordinate \(s_p\), convert them to the same surface material mixture
\(\mathbf f_\Gamma(s_p)\) used by the Stage-C normalization profiles, and
evaluate
\[
q_p=\bar g(\phi_{\rm eff,p},\mathbf f_\Gamma(s_p))
\sqrt{2\bar G(\mathbf f_\Gamma(s_p))H_p},
\qquad
\phi_{\rm eff,p}=\max(\phi_p,0).
\]
Do not use the particle-local composition for either \(\bar G\) or \(\bar g\).
The same \(10^{-4}\) empirical bounded-negative guard used by \(I_h\) applies;
values above one are not clipped.

Under the profile-uniform cohesive assumption, \(q\) is constant through each
normal profile. Use the cached particle-domain-volume weighted consistent Q1
projection to collapse particle \(q\) values to nodal \(T^{\rm coh}_0\).
Commit this field and the simultaneously computed current \(I_{h,0}\) as one
validated initialization operation. Report volume-weighted RMS and maximum
absolute projection residuals, and their normalization by the maximum
profile \(\lvert q\rvert\), as diagnostics. No rejection threshold is inferred.

\subsection{Stage-D lifecycle boundary}

Stage D supplies pure response evaluation and an explicit commit operation for
tests and future solver integration. It does not connect timestep or nonlinear
solver signals. Committed history remains frozen until a later coupled
slip-rate stage has accepted a mechanical timestep.

\section{Particle-to-fault projection}

\subsection{Purpose}
Particle values such as stress-like quantities may need to be represented as Q1 fields on the reconstructed fault.  Particles remain distributed, while every rank owns the complete fault.

\subsection{Geometric projection cache}
For every locally owned particle, cache its optional association with one
finite segment-normal influence strip. The strip half-width is Q1-interpolated
from manager-owned per-vertex projection metadata. For \texttt{dim=2}, a
cache entry should conceptually contain
\begin{lstlisting}[language=C++]
struct ParticleFaultProjection
{
  ParticleId particle_id;
  unsigned int fault_index;
  unsigned int segment_index;
  double xi;
  double particle_domain_volume;
  bool active;
};
\end{lstlisting}

The exact particle-ID type shall follow ASPECT's existing particle infrastructure.

The cache shall be rebuilt when the fault geometry changes or when particle positions/ownership change sufficiently to invalidate the projections.

\subsection{Weights and AMR correction}
Particles are samples, not automatically equal-volume quadrature points.  The existing particle-domain module already computes the physical measure of every particle domain.  For a locally owned particle $p$, obtain
\[
\boxed{
m_p
=
\texttt{particle\_domain\_handler.get\_particle\_domain(}
\texttt{particle.get\_local\_index()).volume()}.
}
\]
This particle-domain measure is the authoritative AMR-aware sampling measure for the particle-to-fault projection.  Do not replace it by the approximation $|K|/N_{p,K}$ unless a future configuration explicitly runs without valid particle-domain data.

Within the admitted parent-domain measure, the geometric kernel is one, so
\[
\sum_q w_{pq}=m_p.
\]
No second user-facing influence-distance or kernel parameter is introduced.
Dynamic phase-field or constitutive values shall not enter this cached weight
unless the cache invalidation rules and authoritative design are first revised.

The CPDI weighting functions provided by \texttt{ParticleDomainAccessor} remain bulk-grid transfer functions.  They shall not be substituted for the reconstructed-fault Q1 shape functions in the least-squares surface projection.

\subsection{Consistent weighted least-squares projection}
\paragraph{Approved discrete revision (2026-09-09).}
The domain-integrated rule replaces the former point-volume formula
and its surface residual/Jacobian counterparts throughout this specification.
The former source obeyed the former specification. Preserve the admitted
parent-center set and each full existing domain $D_p$, without normal clipping.
The objective is now $\sum_p\int_{D_p}(z_p-z_\Gamma(s(x)))^2\,dx$.
Replace $\sum_p m_pN_i(\xi_p)F_p$ by
$\sum_{p,q}w_{pq}N_i(\xi_{pq})F_{pq}$, with the corresponding second Q1
factor in mass and derivative terms; $\sum_qw_{pq}=m_p$.
Bulk FE inputs and particle histories remain parent-center/P0 inputs;
surface Q1 fields and nonlinear responses are evaluated at every domain
quadrature coordinate. This applies consistently to projections, $R_\Gamma$,
$K_V$, $G$, mass-based residual norms and diagnostics.
The first integrated geometry supports straight ordered 2-D faults, including
inclined and nonuniform collinear segments. Split domains at node planes and
projected polygon vertices. Beyond true tips use constant endpoint continuation
without admitting new parents or joining endpoint DoFs. Three-point Gauss
integrates the linear transverse width times Q1 products exactly; nonlinear
accuracy must be tested separately. Curved domain partitions, overlap and 3-D
were unsupported by the first implementation. The approved extension in
\texttt{benchmarking/stage\_K2\_polyline\_quadrature\_addendum.md} defines
the open 2-D polyline map: nearest finite segment-normal projection without
reapplying width, with nearest-vertex continuation in corner/tip gaps. Split
at endpoint planes and distance bisectors and use the selected segment's
one-sided tangent/normal frame. Parent admission/fault assignment, full measure
and parent-P0 evaluation remain unchanged. Overlap, branching, closed loops,
self-intersections and 3-D remain unsupported. History-commit timing is unchanged. The implementation
addendum in \texttt{benchmarking/stage\_K2\_domain\_quadrature\_addendum.md}
defines evaluation locations.

Let $z_p$ be a scalar particle quantity.  Let $N_a(\boldsymbol\xi_p)$ denote the Q1 shape functions of the projected fault cell.  Determine the fault nodal field $\mathbf{z}_\Gamma$ by
\[
\boxed{
\min_{\mathbf{z}_\Gamma}
\sum_p\int_{D_p}
\left(
 z_p-\sum_aN_a(\boldsymbol\xi(\boldsymbol x))z_{\Gamma,a}
\right)^2\,d\boldsymbol x.
}
\]
The normal equations are
\[
\boxed{
\mathbf{M}_\Gamma\mathbf{z}_\Gamma=\mathbf{b}_\Gamma,
}
\]
with
\[
(\mathbf{M}_\Gamma)_{ab}
=
\sum_{p,q} w_{pq}N_a(\boldsymbol\xi_{pq})N_b(\boldsymbol\xi_{pq}),
\]
\[
(\mathbf{b}_\Gamma)_a
=
\sum_{p,q} w_{pq}N_a(\boldsymbol\xi_{pq})z_p.
\]
The same matrix is reused for all projected scalar components that use the same particle set and geometric weights.

For a derived scalar that is not itself a particle property, the manager may
accept finite caller-computed values keyed by stable locally owned particle ID
and return the replicated nodal projection without storing constitutive
meaning. This path shall reuse the same cached mass operator and MPI reduction.
Active associated particles require one value; inactive particles may be
omitted. Volume-weighted RMS and maximum residuals against the projected Q1
field shall be returned as generic diagnostics.

The manager shall also support the reverse local operation needed to evaluate
surface constitutive inputs at particles: Q1-interpolate a registered fault
property at every active locally owned particle's cached fault coordinate and
return the values keyed by stable particle ID. This operation performs no new
geometric search and no MPI communication.

\subsection{MPI assembly}
Each rank assembles only contributions from its locally owned particles:
\[
\mathbf{M}_{\Gamma,r},\qquad
\mathbf{b}_{\Gamma,r}.
\]
Then
\[
\boxed{
\mathbf{M}_\Gamma=\sum_r\mathbf{M}_{\Gamma,r},
\qquad
\mathbf{b}_\Gamma=\sum_r\mathbf{b}_{\Gamma,r}.
}
\]
The sums shall be formed by MPI collective reduction.  No particle positions or particle values are gathered to a coordinator.

All ranks use the same ordered fault vertices and identical tridiagonal Q1
layout. Matrix entries may therefore be packed in deterministic
fault-major/vertex-major order for collective reduction.

\subsection{Cache the matrix, not the nonlinear values}
As long as particle positions, the fault geometry, and the geometric weights remain unchanged,
\[
\mathbf{M}_\Gamma
\]
is constant throughout all nonlinear iterations.

Therefore:
\begin{enumerate}
  \item build the particle projection cache once;
  \item assemble and MPI-reduce $\mathbf{M}_\Gamma$ once;
  \item build a reusable solver/preconditioner/factorization once;
  \item in every nonlinear iteration, assemble only the changing right-hand side $\mathbf{b}_\Gamma$;
  \item MPI-reduce the right-hand side;
  \item solve the small replicated fault system on every rank.
\end{enumerate}

When several particle quantities are projected in the same iteration, their right-hand sides shall be packed into one or a small number of collective reductions.

A coverage diagnostic shall verify that every active fault DoF has sufficient weighted particle support.  Insufficient coverage is an error during the initial implementation rather than a reason to silently invent missing data.

\section{Fault-to-quadrature-point interpolation}

Fault-to-QP interpolation is in the nonlinear hot loop and must be communication-free.

For every relevant local Stokes quadrature point, build a cache containing its closest fault cell, fault reference coordinate, Q1 shape values, signed distance, and geometry version.  For \texttt{dim=2}, conceptually:
\begin{lstlisting}[language=C++]
struct QPFaultProjection
{
  unsigned int fault_index;
  unsigned int segment_index;
  double xi;
  double shape_0;
  double shape_1;
  double signed_normal_distance;
  dealii::Tensor<1,dim> tangent;
  dealii::Tensor<1,dim> normal;
  std::uint64_t geometry_version;
};
\end{lstlisting}

If a fault surface field has nodal values $z_a$, evaluate
\[
\boxed{
z_q=\sum_aN_a(\boldsymbol\xi_q)z_a.
}
\]
For a 2-D segment,
\[
z_q=(1-\xi_q)z_i+\xi_qz_{i+1}.
\]

Because every rank stores the complete surface field, this operation requires no MPI communication.  During nonlinear iterations there shall be no R-tree search, no point-location search, and no fault-matrix rebuild for cached QPs.

Only QPs inside the configured fault influence band need full projection entries; all others may be marked inactive.

\section{Spatial-search strategy}

The ordered-polyline registry provides current 2-D connectivity but is not by
itself an optimal nearest-segment search structure.

The implementation should reuse deal.II \texttt{GridTools::Cache} and related geometry utilities where they directly support the required operation.  If an explicit nearest-fault-cell acceleration structure is still needed, build an application-owned R-tree/AABB index from fault-cell bounding boxes.

For the current 2-D implementation, the final closest-point projection onto a candidate segment is analytic.  A future 3-D closest-point implementation may use the triangulation neighbor structure and standard mapping machinery, but it is outside the present scope.

All spatial-search caches are reconstructible and are invalidated when the fault geometry version changes.

\section{Cache lifetimes}

Three lifetimes shall be distinguished explicitly.

\subsection{Geometry-lifetime caches}
Valid until persistent fault geometry changes:
\begin{itemize}
  \item ordered-segment adjacency and any future acceleration structure;
  \item fault-cell bounding boxes and spatial index;
  \item cell tangents/normals where cached.
\end{itemize}

\subsection{Mesh/particle-lifetime coupling caches}
Valid while the relevant bulk geometry remains unchanged relative to the fault:
\begin{itemize}
  \item particle-to-fault projection entries;
  \item QP-to-fault projection entries;
  \item particle-to-fault projection matrix $\mathbf{M}_\Gamma$;
  \item reusable solver/preconditioner/factorization for $\mathbf{M}_\Gamma$.
\end{itemize}

They are invalidated by events such as fault growth, AMR/repartitioning that changes QPs, or particle movement/migration that changes particle projections.

The adaptive $I_h$ evaluator may retain the geometric point-location maps of
the last preparation's batch sequence. Reuse requires exact batch-coordinate
equality and a valid deal.II mesh lookup on every rank; a mismatch on any rank
requires collective rebuilding. It never caches phase-field values or
constitutive responses. Unused trailing batches are discarded. Configurations
with mesh deformation use fresh lookups, because mapping motion need not emit
a triangulation-change signal. Adaptive panels, tolerances and tail criteria
are unchanged by this geometric reuse.

Cold lookup may reject requests outside a conservative global enclosure of
the reference-tolerance-expanded mesh. The implemented fast path is restricted
to exact MappingCartesian or degree-one MappingQ/MappingQ1 maps on axis-aligned affine hypercubes;
all other mappings/cell shapes retain the original search. Each cell enclosure
includes the unchanged deal.II reference-cell tolerance and an outward
floating-point guard. Bounds are reduced across all owners before rejection.
Surviving requests retain their order and the original distributed ownership
search; samples are scattered back to the original adaptive-request indices.
Rejected entries remain missing, not zero-valued found samples. The enclosure
has the same mesh/mapping lifetime as the geometric lookup cache.

\subsection{Nonlinear-iteration data}
These values are never treated as geometric caches:
\[
\tau,\qquad\sigma_n,\qquad V,\qquad
\text{other current constitutive quantities}.
\]
They are propagated through the cached operators every nonlinear iteration.

\section{Nonlinear-iteration pipeline}

With geometry and projection caches already valid, the high-frequency nonlinear loop is
\[
\boxed{
(\mathbf{x},\mathbf{V}_\Gamma)
\xrightarrow{\;\text{bulk evaluation and cached fault interpolation}\;}
\{F_p,(K_V)_p\}
\xrightarrow{\;\text{local Q1 assembly and MPI reduction}\;}
(\mathbf{R}_\Gamma,\mathbf{K}_V).
}
\]

The intended hot-loop operations are:
\begin{enumerate}
  \item evaluate current bulk trial fields at locally owned associated
        particle points;
  \item interpolate the explicit replicated trial slip-rate vector through
        cached particle/fault coordinates;
  \item evaluate non-committing pointwise constitutive responses;
  \item locally assemble fault-sized Q1 residual and Jacobian coefficients;
  \item perform one packed MPI reduction and factor each replicated $K_V$
        block when a new linearization is requested.
\end{enumerate}

Current trial stress shall not be projected to a frozen fault field because
that would discard its dependence on the Stokes unknowns.  There shall be no
repeated nearest-fault-cell search or particle-projection-matrix assembly
inside the nonlinear iteration.

\section{Optional Xu-style initialization}

If the initial sharp fault is prescribed, Xu's global regression is not used.

If a future application requires reconstruction of an initial sharp fault from a phase-field point cloud, a weighted form of Xu et al.'s Eq. (18) may be used:
\[
(\mathbf{A}^T\mathbf{W}\mathbf{A}+\lambda\mathbf{D}^T\mathbf{D})\boldsymbol\beta
=
\mathbf{A}^T\mathbf{W}\mathbf{Y}.
\]

Even if millions of bulk nodes are active, the point cloud shall not be gathered.  Each MPI rank processes only uniquely owned FE nodes and assembles local sufficient statistics
\[
\mathbf{H}_r
=
\sum_{i\in\mathcal P_r}w_i\mathbf{A}_i^T\mathbf{A}_i,
\qquad
\mathbf{g}_r
=
\sum_{i\in\mathcal P_r}w_i\mathbf{A}_i^TY_i.
\]
Then
\[
\mathbf{H}=\sum_r\mathbf{H}_r+\lambda\mathbf{D}^T\mathbf{D},
\qquad
\mathbf{g}=\sum_r\mathbf{g}_r
\]
are formed by MPI reduction and the reduced fault problem is solved on the replicated fault.

For \texttt{dim=2}, each interpolation row $\mathbf{A}_i$ has only two nonzero entries, so the reduced matrix is sparse/banded.  A bulk-node-to-fault cache is unnecessary because this initialization is performed rarely and is not in the nonlinear hot loop.

\section{AMR, particle migration, and cache invalidation}

The persistent fault registry is independent of the bulk triangulation.  Bulk AMR or repartitioning does not transfer or modify fault state.

After AMR/repartitioning:
\begin{itemize}
  \item persistent fault geometry/state remain unchanged;
  \item distributed phase-field point-evaluation structures are rebuilt as required;
  \item QP-to-fault caches are invalidated because the local cells/QPs may have changed;
  \item particle-to-fault caches are rebuilt if particle positions/ownership changed;
  \item fault geometry-dependent search/connectivity caches are rebuilt only
        if fault geometry changed.
\end{itemize}

\section{Geometry versioning}
Maintain
\begin{lstlisting}[language=C++]
std::uint64_t geometry_version;
\end{lstlisting}
and increment it whenever new persistent geometry is committed.

Every geometry-dependent cache stores the version from which it was built.  A version mismatch makes the cache invalid.

Additional bulk-mesh and particle-position/update counters may be used for QP and particle caches.

\section{Checkpoint/restart}

Checkpoint persistent data:
\begin{itemize}
  \item ordered fault vertices and implicit segment connectivity;
  \item projection half-width metadata;
  \item the generic fault-property schema and committed values;
  \item timestep-committed slip rate;
  \item fault constitutive state;
  \item geometry version.
\end{itemize}

Do not checkpoint reconstructible runtime data:
\begin{itemize}
  \item triangulation/DoF numbering caches if they can be rebuilt from the registry;
  \item R-trees/AABB structures;
  \item tangents/normals;
  \item particle/QP projection caches;
  \item distributed point-evaluation caches;
  \item factorization/preconditioner objects.
\end{itemize}

After restart, rebuild property-name indices, current slip-rate state, and all
projection/search caches from the serialized registry.

\section{Timestep ordering}

For the initial implementation use
\[
\boxed{
\text{solve bulk fields}
\rightarrow
\text{update phase field}
\rightarrow
\text{reconstruct/extend fault}
\rightarrow
\text{initialize state on new fault vertices}.
}
\]

New persistent geometry becomes active in the constitutive calculation of the next timestep.  Fault topology shall not change in the middle of a nonlinear Stokes solve.

Within a nonlinear solve, the fault geometry is therefore fixed and all geometry-dependent coupling caches can be reused.

\section{Coupled bulk--surface solver architecture}

The coupled implementation shall preserve a strict responsibility boundary.
\texttt{MaterialModel::PhaseFieldFault} directly exposes the narrow pointwise
input, response, and semantic operations needed by reconstructed-fault
coupling; no abstract reconstructed-fault constitutive base class is used.  It
shall not assemble finite element weak forms and shall not expose
\texttt{apply\_B()} or \texttt{apply\_G()} through the material-model
interface.  A simulator-owned computational helper named
\texttt{ReconstructedFaultSurfaceSystem} checks for the concrete
\texttt{PhaseFieldFault} material model once at construction and retains that
reference for its lifetime.  The simulator owns one canonical surface-system
instance and one canonical reconstructed-fault Stokes-coupling instance when
the coupled model is available.  Solver objects retain references to these
components and shall not construct duplicate stateful instances.  The surface
helper owns the particle/Q1 surface assembly,
replicated fault-vector MPI reductions, and the surface factorization.  This
particle workflow is not an ordinary ASPECT cell-local Scratch/CopyData
assembler.  The condensed operator remains solver-owned.

\subsection{Non-committing point response}

At an associated locally owned particle, let
\begin{equation}
 F=t-T^{\rm coh}-\mu(V,\text{state})\sigma_n-\eta^d V,
 \qquad
 t=\boldsymbol\tau:\boldsymbol S.
\end{equation}
The pointwise material capability evaluates $F$, its negative derivative with
respect to $V$, and the constitutive coefficients required by the later $B$
and $G$ actions.  The input carries the explicit bulk and slip-rate trial
state.  Evaluation shall not change the manager's current or committed slip
rate, committed fault properties, particle Maxwell stress, or any other
history.  This same non-committing path is used for Newton linearization and
line-search residual evaluation.

At timestep zero, the positive parsed \texttt{Initial time step} supplies the
Maxwell interval; after timestep zero, use the current simulator timestep.
Surface evaluation requires current $I_h$, committed cohesive history, and
initialized slip rate.  Stateful rate-and-state friction additionally
requires a finite positive committed $\Theta$ at every fault vertex.  Stage F
registers and validates that state but does not invent its initial value or
commit rule; the coupled lifecycle in Stage I must define and connect those
operations before the production solve may use rate-and-state friction.

\subsection{Normal-pressure modes}

\texttt{PhaseFieldFault}, not \texttt{FaultFriction}, owns the boolean
parameter \texttt{Use adiabatic pressure in fault friction}, whose default is
false.  The two modes are
\begin{equation}
 \begin{aligned}
  \text{false:}\quad &\sigma_n=p-\boldsymbol\tau:\boldsymbol N,\\
  \text{true:}\quad  &\sigma_n=p_{\rm ad}(\boldsymbol x).
 \end{aligned}
\end{equation}
In adiabatic-pressure mode the adiabatic model pressure is the complete normal
pressure in the friction term.  Dynamic pressure and deviatoric normal
traction are not added to it.  An enabled adiabatic-pressure mode requires
initialized adiabatic conditions.

\subsection{Stage-F surface residual and Jacobian}

The replicated surface residual has one nested vector per reconstructed fault
and one entry per Q1 fault vertex.  Only locally owned associated particles
assemble
\begin{equation}
 (R_\Gamma)_i=\sum_{p,q} w_{pq}N_i(\xi_{pq})F_{pq}.
\end{equation}
Fault-sized local vectors are summed over MPI so every rank obtains the same
replicated residual.  Report volume-weighted RMS residuals per fault and over
all faults; these are diagnostics rather than additional equations.

With fixed geometry, phase field, particle history, and committed $\Theta$,
define
\begin{equation}
 K_V=-\frac{\partial R_\Gamma}{\partial V}.
\end{equation}
For the current constitutive law its Q1 entries are
\begin{equation}
 (K_V)_{ij}
 =\sum_{p,q} w_{pq}N_i(\xi_{pq})N_j(\xi_{pq})
 \left[
  2\kappa_p\chi_p\boldsymbol S_p:\boldsymbol S_p
  +\frac{\kappa_p}{I_{h,p}}
  +\sigma_{n,p}
    \left.\frac{\partial\mu}{\partial V}\right|_{\Theta}
  +\eta^d_p
 \right],
 \qquad \chi_p=\frac{h_p}{I_{h,p}}.
\end{equation}
The rate-and-state derivative is the exact partial derivative at fixed
$\Theta$; the rate-dependent law supplies its exact signed derivative.  The
current ordered Q1 polyline and single-segment association make every fault
block symmetric tridiagonal.  The signed friction contribution means that
$K_V$ is not assumed positive definite.

Store the reduced diagonal and off-diagonal coefficients as replicated
per-fault data.  For the initial implementation, materialize each block as a
deal.II sparse matrix and factor it with the tested
\texttt{SparseDirectUMFPACK} implementation, which supports nonsingular
indefinite matrices.  Do not add a bespoke tridiagonal factorization or
condition estimator.  Factorization failure is a numerical failure.  Every
application of the stored inverse shall also check a scaled backward error and
diagnose a singular or excessively ill-conditioned block explicitly.

The surface-system helper exposes a non-committing residual evaluation, an
operation that assembles/replaces the current $K_V$ linearization and factors,
and an operation that applies the stored replicated $K_V^{-1}$ to a
fault-major vector.  This is the only inverse interface required by later
condensation.  A caller can evaluate arbitrary trial $V$ values through the
explicit input without changing manager-owned nonlinear state.

Starting a new linearization invalidates the previous one.  Assemble and
factor a complete candidate and publish it only after every fault block
succeeds.  A failed factorization shall therefore leave no stale or partially
factorized $K_V$ available.  The helper is a simulator-owned, non-checkpointed
computational component.  Its surface linearization is valid only as part of
one coupled linearization and introduces no constitutive state or additional
timestep lifecycle hook.

Stage F ends at this surface system.  It shall not assemble the
slip-dependent bulk residual, implement either coupling action, modify the
Stokes operator, connect a nonlinear solver lifecycle, or commit history.

\subsection{Stage-G weak-form coupling}

Stage G introduces \texttt{Assemblers::ReconstructedFaultStokes} as a genuine
cell/QP \texttt{Assemblers::Interface} implementation.  It owns the
slip-dependent bulk Stokes residual and the fault-to-bulk $B$ action.  The
particle-based bulk-to-surface $G$ action instead extends
\texttt{ReconstructedFaultSurfaceSystem}, where it reuses the Stage-F point
response values $\kappa$, $\chi$, and $\mu$, particle associations, and the
replicated fault-vector reduction.  The solver orchestrates these independent
actions.  The fault-induced bulk residual is
\begin{equation}
 R_{\rm fault}(\boldsymbol w;V)
 =-\int_\Omega 2\kappa(\chi V+\upsilon^{\rm hist})
   \boldsymbol S:\dot{\boldsymbol\epsilon}(\boldsymbol w)\,d\Omega.
\end{equation}
Consequently the standalone residual operation overwrites its output with
$R_{\rm fault}$, the standalone $B$ operation overwrites its output with
$+B\,\delta V$, and the ordinary assembler adds $-R_{\rm fault}$ into the
Stokes right-hand side.  The frozen
$-2\kappa\upsilon^{\rm hist}\boldsymbol S$ contribution is part of the
absolute residual even though it does not contribute to $B$.

The same ordinary cell assembler supplies the frozen Maxwell-history load
\[
 -\int_\Omega \beta\boldsymbol\tau_{\rm old}:
       \dot{\boldsymbol\epsilon}(\boldsymbol w)\,d\Omega
\]
to the Newton right-hand side exactly once, over the entire bulk, not just
fault-associated QPs. Its positive counterpart is part of the bulk residual.
Evaluate $\beta\boldsymbol\tau_{\rm old}$ pointwise using the local bulk
composition and temperature. The explicitly mapped stress compositional
fields in the current working FE vector represent the frozen particle
history; do not read an older FE timestep or evolve that history in a trial.
The slip-only residual and $B$ actions remain unchanged. Timestep-zero history
retention and terminal history publication retain their established semantics.

For particle-to-FE transfer, shared continuous DoFs receive the arithmetic
average of incident-cell interpolator proposals. Accumulate values and counts
from locally owned cells with MPI ADD, count each cell once, and divide at
owned DoFs. Unshared/DG DoFs retain their single proposal. This generic rule
also applies to other particle-mapped continuous fields; it is not an L2
projection or a change of particle interpolator. Apply physical constraints
only through the existing private working-vector workflow, not by republishing
history.

Accumulate the genuinely frozen Maxwell/profile right-hand side separately
from unknown-dependent contributions through local quadrature, homogeneous
constraint distribution, and MPI compression. Combine only the completed
global vectors. The $BV$ contribution remains unknown-dependent, including
when an active set happens to give $\delta V=0$. Use this same evaluation for
every Newton base and non-committing trial. This arithmetic separation avoids
cancellation of small residuals against stress loads without changing the
equations, normalization scales, or convergence and line-search criteria.

A fault-slip increment changes stress by
\begin{equation}
 \delta\boldsymbol\tau
 =-2\kappa\chi\boldsymbol S\,\delta V_\Gamma,
\end{equation}
which defines $B$ after weak-form insertion.

For dynamic fault pressure, the bulk perturbation of the surface equation is
\begin{equation}
 G\delta x
 =2\kappa(\boldsymbol S+\mu\boldsymbol N):
   \delta\dot{\boldsymbol\epsilon}
  -\mu\,\delta p.
\end{equation}
For adiabatic fault pressure, $p_{\rm ad}$ is prescribed and has no variation
with the Stokes unknowns, so
\begin{equation}
 G\delta x
 =2\kappa\boldsymbol S:\delta\dot{\boldsymbol\epsilon}.
\end{equation}
There is no $\mu\boldsymbol N$ contribution and no $-\mu\delta p$ term in
this mode.  The Stokes assembler performs no fault-vector MPI reduction and
the surface helper performs no cell-local bulk weak-form assembly.  Stage G
shall include finite-difference tests of both modes, but it shall not implement
condensation or nonlinear lifecycle.

Maintain two caches with distinct lifetimes.  The manager-owned geometry cache
maps every locally owned cell and every quadrature point of the exact
production Stokes velocity quadrature, in its original order, to the fault,
segment, Q1 coordinate and weights, signed distance, position, tangent, and
normal.  It persists only while mesh and mapping-relevant geometry, fault
geometry, projection widths, and quadrature identity are unchanged.  Relevant
mesh, restart, deformation, and fault/projection mutations invalidate it.
Debug builds shall verify the exact reference points, weights, count, and
physical point ordering at consumer boundaries.  The assembler-owned
linearization cache stores $2\kappa\chi\boldsymbol S$ at each associated QP.
Rebuild it once from an explicit physical bulk linearization state.  Repeated
$B$ applications shall use only this frozen cache and the geometry cache; they
shall not reevaluate the material model or repeat fault reconstruction.  A
geometry-cache change invalidates use of the linearization cache.

ASPECT full solution, old-solution, and nonlinear-linearization vectors store
physical pressure.  The Stokes solver vector uses
\begin{equation}
 \widehat p=\frac{p_{\rm physical}}{s_p},
 \qquad s_p=\texttt{get\_pressure\_scaling()}.
\end{equation}
Stage-G residual and $G$ interfaces accept physical full-system vectors.
Stage H shall convert a solver direction according to
$\delta p_{\rm physical}=s_p\delta\widehat p$, equivalently
$G_{\rm solver}=G_{\rm physical}\operatorname{diag}(I_u,s_pI_p)$.
The dynamic-pressure term for a solver direction is therefore
$-\mu s_p\delta\widehat p$.  Adiabatic-pressure mode has no pressure term.
Do not implement this solver-vector adapter in Stage G.

\subsection{Stage-H exact condensation}

The Newton block equation and its right-hand side are
\begin{equation}
 \boxed{
 \begin{bmatrix}
  A & -B\\
  G & -K_V
 \end{bmatrix}
 \begin{bmatrix}
  \delta x\\
  \delta V
 \end{bmatrix}
 =-
 \begin{bmatrix}
  R_{\rm bulk}\\
  R_\Gamma
 \end{bmatrix}.}
\end{equation}
Therefore the solver-owned condensed equation and recovery are
\begin{equation}
 \boxed{
 (A-BK_V^{-1}G)\delta x
 =-R_{\rm bulk}+BK_V^{-1}R_\Gamma,}
\end{equation}
\begin{equation}
 \boxed{
 \delta V=K_V^{-1}(R_\Gamma+G\delta x).}
\end{equation}
Stage H adds only this exact condensed operator, right-hand side, and recovery.
It references the canonical simulator-owned surface and Stokes-coupling
components and owns no constitutive state.  The assembled $A$, frozen $B$, and
surface $G/K_V$ data form one coupled linearization: their validity begins only
after all components have been built successfully and ends when any component
is rebuilt or invalidated.  This lifetime is independent of the lifetime of
the helper object.

The condensed implementation shall depend only on semantic $B$ and $G$ actions
and a semantic surface linear solve.  Stage H provides an unrestricted
$K_V^{-1}$ implementation.  Stage I may substitute a solve on the active/free
partition, projecting its right-hand side and returning zero active
components, without changing the condensation implementation.

Bulk solver directions are homogeneous perturbations.  Constrained algebraic
entries are zero when applying $A$.  Before evaluating the physical finite
element field for $G$, distribute homogeneous hanging-node and periodic
constraints while keeping every prescribed-boundary inhomogeneity zero.  The
$B$ action uses the same homogeneous constraints, and constrained entries of
the returned algebraic vector are zero.  Inhomogeneous solution values shall
never enter a Jacobian direction.

Before the first coupled residual at each timestep, distribute the current
physical inhomogeneous constraints into a private bulk base iterate. This
includes initial nonzero and later changed prescribed velocities. Subsequent
Newton residual and matrix assembly uses homogeneous Stokes constraints: the
lift is already present in the physical iterate and shall not be subtracted
again by local-to-global elimination. Preserve auxiliary-field constraints and
restore the caller's constraints on exit. A failed solve shall discard this
private lift along with all accepted nonlinear updates, without changing the
production bulk solution or timestep-committed $V$.

For prescribed adiabatic friction pressure, apply the existing ASPECT
\texttt{Pressure normalization} rule to the private physical base and every
physical trial before residual and merit evaluation. The modes
\texttt{volume}, \texttt{surface}, and \texttt{no} retain their standard
mean-pressure semantics. Do not normalize a homogeneous direction to an
inhomogeneous surface-pressure target or change solver pressure scaling.
Both published bulk vectors shall inherit the same normalized accepted
iterate. Its pressure-normalization adjustment remains private until the
terminal commit; rejection or failure leaves published bookkeeping unchanged.
Every normalization allocation and MPI collective precedes all history and
$V$ writes. This gauge-only correction shall not shift pressure in the
true-normal-stress friction mode, where it would change the surface equation.

The supported assembled iterative path shall use FGMRES.  The condensed
operator is generally nonsymmetric because $B$ and $G$ are not assumed to be
transposes, so CG and MINRES are not valid outer solvers.  Verification shall
also check a freshly computed residual of every returned linear direction;
the Arnoldi residual estimate alone is insufficient. Residual replacement and
restart share the existing total iteration budget. In incompressible,
prescribed-friction-pressure closed/periodic configurations only, a verified
constant-pressure null mode may be removed algebraically. Verify both null
identities for the full homogeneous constrained condensed operator; $B$ has
no pressure rows, while the right-null test includes $G$ and the current
surface inverse. Project operator and preconditioner consistently. Reject
RHS or residual incompatibility larger than either $100\epsilon_{\rm mach}$
times the maximum of the initial bulk residual, zero-velocity reference,
and current condensed RHS norm, or the unchanged absolute nonlinear bulk
target. Retain and report raw and projected residuals. This numerical
backward-error check introduces no physical parameter and does not alter
physical pressure normalization, homogeneous Newton constraints, or true-
normal-stress/absolute-pressure configurations. Verification shall
include a small explicit solve of the complete two-block system compared with
the condensed solve and recovery, and a multiple-fault case.  Stage H does not
connect nonlinear commit/rollback.

\subsection{Stage-I lower bound and nonlinear lifecycle}

The coupled solver works directly with $V\geq V_{\min}$.  Before applying a
fraction-to-boundary limit, form a lower-bound active set and projected Newton
direction.  If a degree of freedom is at $V_{\min}$ and its Newton direction
points below the bound, hold that degree of freedom active and solve/project
the remaining free direction consistently.  Apply fraction-to-boundary only
to the free direction.  Fraction-to-boundary alone is insufficient because an
outward direction at an active node would otherwise force
$\alpha_{\max}=0$.

Every line-search candidate uses the non-committing coupled residual path.
Rejected candidates do not accumulate.  Only an accepted trial replaces the
current Newton iterate, and only nonlinear convergence commits timestep state.
A failed solve or rejected timestep restores current $V$ from committed $V$
and leaves committed $\Theta$, $T^{\rm coh}$, previous $I_h$, and particle
Maxwell stress unchanged.

Use the local bound test
\begin{equation}
 V_i-V_{\min}\leq 100\epsilon_{\rm mach}\max(V_{\min},|V_i|).
\end{equation}
For each Newton iteration, start from an empty active set and add an at-bound
vertex only when its recovered direction is negative. The restricted inverse
is the principal free-block solve $K_{\mathcal F\mathcal F}^{-1}$; active rows
and columns are disconnected, the right-hand side is projected to the free
set, and active increments are exactly zero. Recompute the active set from
empty only after an accepted step.

For free vertices with negative direction use the exact fraction-to-boundary
ratio $(V_i-V_{\min})/(-\delta V_i)$, without a strict-interior factor. Project
roundoff-level contacts exactly to $V_{\min}$, but do not clamp constitutive
inputs. Hold the active set fixed during the line search and accept a trial
only if
\begin{equation}
 \Phi_{\rm trial}\leq (1-10^{-4}\alpha)\Phi_k,
 \qquad \Phi=\frac12(r_b^2+r_\Gamma^2).
\end{equation}
Start from the admissible maximum step and reduce by $2/3$. Exhaustion is a
nonlinear failure and shall not accept the last unsuccessful candidate.

Bulk and projected-free surface residuals use separate, fixed normalization
scales during a mechanical solve and must independently satisfy the nonlinear
tolerance. The bulk scale floor retains the initial Stokes residual convention
using $\max(\epsilon_{\rm linear},\sqrt{\epsilon_{\rm mach}})$.
Include the independently calculated absolute bulk precision scale
\[
 d_i=\sum_{c\in\{u,p\}}\sum_{j\in c}|(A_0)_{ij}|\,
                  \|(x_0)_c\|_\infty,
 \qquad \rho_b=\epsilon_{\rm mach}\|d\|_2,
 \qquad x_0=(u_0,p_0/s_p).
\]
Here $A_0$ and $x_0$ are the first homogeneous-constraint bulk matrix and
physical iterate in solver coordinates. Absolute row sums bound the response
to machine-precision bulk perturbations at those block magnitudes; MPI sums
count owned rows once. This is not a fitted stagnation floor or a physical
error estimate. The mixed bulk criterion is
$\|R_b\|<\epsilon_{\rm nl}s_b+\rho_b$. Use
$s_b^{\rm mixed}=s_b+\rho_b/\epsilon_{\rm nl}$ in both convergence and merit,
fixed throughout the solve. Pressure-compatibility projection retains its
original relative-target cap; fresh linear verification and the independent
surface criterion are unchanged. Bulk FE gradients are evaluated after
removing a cellwise constant velocity, whose gradient is exactly zero, to
reduce assembly cancellation. This cannot remove rounding of the represented
iterate. Stagnation above the mixed target remains failure.
The surface scale is fixed at the first stabilized linearization as
\[
s_\Gamma=\max\bigl(\|P_{\mathcal F}R_{\Gamma,0}\|_\Gamma,
                     \|R_{\Gamma,0}\|_\Gamma,
                     \|K_V V_{\rm char}\|_\Gamma\bigr),\qquad
V_{{\rm char},i}=\max(V_{\min},|V_i|),
\]
with both reference norms evaluated on all vertices. This characteristic
traction scale is not suppressed by a roundoff multiplier, and is used
unchanged for both convergence and the Armijo merit throughout the solve.
No dimensional parameter is added; configured tolerances and the line-search
budget are unchanged. If a block and its reference scale are exactly zero,
its normalized residual is zero only while that block remains exactly zero.

The surface residual is a weak Q1 nodal vector. Let $M_\Gamma$ be the
consistent Q1 mass matrix assembled with the same particle-domain quadrature
as $R_\Gamma$. For the current free set, use the strong-residual RMS norm
\begin{equation}
 \left\|P_{\mathcal F}R_\Gamma\right\|_\Gamma
 =\left[
 \frac{(P_{\mathcal F}R_\Gamma)^T
 M_{\mathcal F\mathcal F}^{-1}
 (P_{\mathcal F}R_\Gamma)}
 {{\boldsymbol 1}_{\mathcal F}^T
 M_{\mathcal F\mathcal F}{\boldsymbol 1}_{\mathcal F}}
 \right]^{1/2}.
\end{equation}
Both the residual and mass matrix are restricted to the same active/free
partition. Define the norm as zero when all vertices are active. The mass
matrix is geometry/quadrature data and remains fixed for one coupled
linearization.

Mechanical preparation may initialize missing state only in a fresh
timestep-zero model. It recomputes current $I_h$, initializes missing cohesive
history, sets $V=V_{\min}$, and obtains a positive $\Theta_0$ by projecting
exactly one generic particle-advected compositional field mapped to component
zero of particle property \texttt{phase field fault state}. Restarted or
later-time models require complete committed state. Stage I commits only $V$.

\subsection{Stage-J history feedback}

The authoritative fixed-fault ordering is
\begin{equation}
 H_{k-1}\longrightarrow\phi_k\longrightarrow(u_k,p_k,V_k)
 \longrightarrow\{\Theta_k,T_k^{\rm coh},\tau_k,H_k,I_{h,k}\}
 \longrightarrow\phi_{k+1}.
\end{equation}
The phase field, reconstructed geometry, and all committed histories remain
fixed throughout one coupled mechanical solve. All candidate construction,
MPI communication, allocation, and validation shall finish before a terminal
non-throwing mutation phase commits histories, $V$, and the accepted bulk
solution.

Timestep zero commits the solved initial kinematic state $V_0$ but does not
physically evolve $\Theta_0$ or $H_0$ with the artificial
\texttt{Initial time step}. It retains the supplied $\Theta_0$, initialized
irreversible $H_0$, initialization-specific $T_0^{\rm coh}$, and initialized
particle Maxwell stress. It shall store $I_{h,0}$ as the previous-normalization
snapshot required by the first real timestep.

For the initial mechanical evaluation, use the converged $\phi_0$ as both
current and previous phase profile, consistently with this initialized
previous-$I_h$ snapshot. The initially zero phase block in
\texttt{old\_solution} is not an earlier physical profile. The shared pointwise
localization evaluation used by surface and bulk coupling shall apply this
rule without overwriting the old finite-element vector or advancing any
retained history. At $k>0$, including after restart, continue to use the actual
previous finite-element phase field.

For $k>0$, distinguish bulk and surface Maxwell coefficients. Bulk
coefficients use particle-local material fractions and temperature. Surface
cohesive coefficients use the projected Q1 surface mixture and a transient Q1
fault-surface temperature obtained by sampling the frozen FE temperature at
fault vertices. The surface coefficients govern cohesive traction, cohesive
history correction, and $\partial T^{\rm coh}/\partial V=\kappa_\Gamma/I_h$;
bulk coefficients govern Maxwell stress and the stress parts of $B$ and $G$.

Consistently project accepted particle samples of $T_k^{\rm coh}$ to the
replicated fault Q1 space. Interpolate this projected traction back to each
associated particle before computing its driving force. For $g_k<1$, use the
algebraically exact form
\begin{equation}
 a=\frac{T_k^{\rm coh}}{g_k},\qquad
 b=\frac{\beta_{\Gamma,k}h_{k-1}T_{k-1}^{\rm coh}}{1-g_k},\qquad
 \mathcal H_k=\frac{\Delta t_k}{2\kappa_{\Gamma,k}}(a-b)(a+b),
 \qquad H_k=\max(H_{k-1},\mathcal H_k).
\end{equation}
A negative $\mathcal H_k$ shall not be independently clipped. For the intact
case $g_k=1$, $h_k=h_{k-1}=0$, use
\begin{equation}
 \mathcal H_k=\frac{\Delta t_k}{2\kappa_{\Gamma,k}}(T_k^{\rm coh})^2.
\end{equation}
The case $g_k=1$ with $h_{k-1}>0$ is inadmissible healing.

Rate-and-state $\Theta$ is updated at the fault Q1 vertices with the accepted
$V_k$ through the existing exact aging law and the same projected surface
mixture used by friction. The present law has one global $D_c$; no duplicate
or composition-dependent parameter is added. Rate-dependent friction has no
state update.

Expose the existing law-specific RSF restriction through an opt-in ASPECT
time-stepping plugin named \texttt{reconstructed fault time step}. It uses
timestep-committed Q1 $V$, Q1 surface fractions, the global ASPECT CFL number,
and \texttt{use\_operator\_splitting=true}. Standard ASPECT model-list
selection is unchanged: the plugin is never activated implicitly. Stage J
does not add a post-solve cutback or repeat operation.

\subsection{Mandatory block-Jacobian verification}

Before the coupled API is considered complete, evaluate a single
non-committing coupled residual
\begin{equation}
 R(x,V)=\begin{bmatrix}R_{\rm bulk}(x,V)\\R_\Gamma(x,V)\end{bmatrix}
\end{equation}
and compare centered finite differences with the uncondensed block action.
Tests shall cover velocity-only, pressure-only, slip-only, and mixed
directions; both normal-pressure modes; nonzero cohesive and evolving-profile
history; one and two MPI ranks; and states near but not crossing $V_{\min}$.
Step-size sweeps shall exhibit the expected truncation-error regime followed
by roundoff saturation.  Independently verify the condensed action, condensed
right-hand-side sign, and recovered fault row.

\section{Diagnostics and assertions}

Required reconstruction diagnostics include
\[
\rho,\qquad d_{\max},\qquad\delta,\qquad\sigma,
\]
and for a 2-D propagating tip,
\[
A_T,\qquad\Delta\theta.
\]

Required coupling diagnostics include
\begin{itemize}
  \item particle support/coverage per active fault DoF;
  \item minimum and maximum diagonal entries of $\mathbf{M}_\Gamma$;
  \item solver convergence for particle-to-fault projections;
  \item number of cached relevant particles and QPs;
  \item cache rebuild counters;
  \item MPI consistency checks of replicated surface vectors.
\end{itemize}

After every fault synchronization, all ranks shall have identical persistent geometry, connectivity, and surface state.

\section{Initial parameter scales}

Initial reconstruction scales (with $\phi_{\min}$ and $\phi_{\max}$ obtained from the existing phase-field material-model interface):
\[
R_n=\max(2\ell,2h_{\mathrm{pf}}),
\qquad
\Delta\eta\approx0.5h_{\mathrm{pf}},
\qquad
h_\Gamma\approx0.5h_{\mathrm{pf}},
\]
\[
p=2,
\qquad
R_T\approx2\ell\text{--}3\ell,
\qquad
R_N\approx2\ell,
\]
\[
\theta_{\max}\approx45^\circ,
\qquad
N_{\mathrm{geom,max}}=3.
\]

The following remain configurable and require calibration:
\[
\rho_{\mathrm{tip}},\quad
d_{\mathrm{tip}},\quad
\lambda_s,\quad
A_0,\quad
\sigma_{\max},\quad
\delta_{\max},
\]
for future propagation geometry. The current particle influence half-widths
come from stationary-profile support and the admitted-strip kernel is one.

\section{Implementation stages}

\subsection*{Stage A: ownership and naming}
Establish the generic manager/material ownership split and generic friction-law
name without changing mechanics.

\subsection*{Stage B: Maxwell stress history}
Add the particle Maxwell-stress property and private non-rotational
time-discrete Maxwell operations. Keep particle history frozen until a future
post-convergence update.

\subsection*{Stage C: distributed adaptive \(I_h\)}
Compute transient current \(I_h\) from distributed Q1 phase field, fixed
surface mixtures, adaptive two-sided profiles, and consistent Q1 projection.

\subsection*{Stage D: common cohesive state}
Register committed \(T^{\rm coh}\) and previous \(I_h\), initialize them from
the prescribed initial \(H\), and implement the exact non-committing cohesive
and history-corrected \(\upsilon\) operations.

\subsection*{Stage E: generic fault friction}
Preserve rate-and-state friction and add the approved rate-dependent law.

\subsection*{Stage F: surface residual and \(K_V\)}
Add only the pointwise non-committing constitutive capability, replicated
surface-system helper, indefinite-capable \(K_V\) factorization, and inverse
application described above.

\subsection*{Stage G: reconstructed-fault Stokes coupling}
Add the dedicated cell/QP assembler's slip-dependent bulk residual and
fault-to-bulk $B$ action, and add the particle-based $G$ action to the surface
helper, including both pressure modes.

\subsection*{Stage H: exact condensed solve}
Add the solver-side \(A-BK_V^{-1}G\) action, the correctly signed condensed
right-hand side, and recovery of \(\delta V\).

\subsection*{Stage I: coupled nonlinear lifecycle}
Add lower-bound active-set/projected directions before fraction-to-boundary,
line search, separate residual convergence, and accepted-state
commit/rollback.

\subsection*{Stage J: history feedback}
Atomically commit the accepted cohesive, rate-and-state, Maxwell, and
crack-driving histories using the ordering and timestep-zero semantics above.

\subsection*{Stage K: fixed-fault benchmarks}
Verify fixed-fault benchmarks. Propagation, 3-D surfaces, and overlapping
faults remain deferred.

\section{Rules for Codex}

When implementing this specification, Codex shall:
\begin{enumerate}
  \item treat this document as authoritative for scientific and architectural decisions;
  \item implement only the requested stage;
  \item use the replicated ordered-polyline registry and generic property
        storage specified here for the current 2-D implementation;
  \item not add a duplicate runtime fault mesh without first revising the
        authoritative documents;
  \item not store persistent state in triangulation user data/flags;
  \item reuse existing deal.II/ASPECT geometry, connectivity, FE, and MPI infrastructure whenever it satisfies the stated requirement;
  \item obey Section~\ref{sec:existing-interface}: reuse the authoritative phase-field, particle-domain, and RSF APIs and do not make the reconstructed-fault code depend on the old core-phase-field/phase-field-RSF architecture;
  \item use application-owned caches for hot-loop projections and interpolation weights;
  \item not gather distributed bulk nodes or particles when a local-assembly plus MPI-reduction formulation exists;
  \item not change formulas, parameter meanings, ownership rules, or cache lifetimes silently;
  \item stop and report any implementation-critical ambiguity or conflict with existing ASPECT/deal.II code before inventing a new scientific or architectural choice;
  \item after each stage, list changed files and map each implementation change back to the corresponding requirement in this document.
\end{enumerate}

\section*{Reference}
Y. Xu, T. You, and Q. Zhu, ``Reconstruct lower-dimensional crack paths from phase-field point cloud,'' \emph{International Journal for Numerical Methods in Engineering}, 124 (2023), 3329--3351, DOI: 10.1002/nme.7249.

\end{document}
